\section{Characterization of the Birkhoff polytope}
In this section, we prove Theorem~\ref{main:birkhoff_unique}.
We first show the following weaker result.

\begin{lemma}\label{l:snbirk_unique}
  Let $D\colon S_n \to \GL(d,\reals)$ be a representation 
  such that the representation polytope~$P(D)$ is combinatorially
  equivalent to the Birkhoff polytope.
  Then $D$ is effectively equivalent to 
  the standard permutation representation $P$ of $S_n$.  
\end{lemma}

\begin{proof}  
  We have to show that $D$ has the same nontrivial constituents
  as $P$, up to automorphisms of $S_n$.
  Since we can replace $D$ by a stably equivalent representation,
  we may (and do) assume that
  the trivial character is not a constituent of the character of $D$.
  
  A combinatorial isomorphism from the Birkhoff polytope~$B_n$
  onto $P(D)$ sends a vertex $P(g)$ of $B_n$ (where $g\in S_n$)
  to a vertex $D(\alpha(g))$ of $P(D)$,
  where $\alpha\colon S_n \to S_n$ is a permutation of $S_n$.
  Then the map sending
  $\gamma\in \Sym{S_n}$ to 
  $\alpha \circ \gamma\circ \alpha^{-1}$
  is an isomorphism from the combinatorial symmetry group of 
  $B_n$ onto the combinatorial symmetry group of $P(D)$.
  The combinatorial symmetry group of the Birkhoff
  polytope is $\Gamma(S_n)$,
  and the combinatorial symmetry group of $P(D)$ contains 
  $\Gamma(S_n)$ (in its natural action on $P(D)$), 
  by Lemma~\ref{l:reppoltrivsyms}.
  Therefore, the combinatorial symmetry group of $P(D)$
  is just $\Gamma(S_n)$.
  It follows that $\alpha \in \Norml_{\Sym{S_n}}(\Gamma(S_n))$.
  By Lemma~\ref{l:sn_cent_est}, Corollary~\ref{c:normlgamma}
  applies to $S_n$ and thus 
  $\alpha\in (\Aut S_n)\Gamma(S_n)$.
  After multiplying $\alpha$ with an element of
  $\Gamma(S_n)$, we may thus assume
  $\alpha \in \Aut S_n$.
  Since then $D$ and $D\circ\alpha$ are effectively equivalent,
  we may assume that $\alpha = \Id_{S_n}$.
  This means that the combinatorial isomorphism from $B_n$ onto $P(D)$
  simply sends the vertex $P(g)$ to $D(g)$, for any $g\in S_n$.
  In particular, a subset of $S_n$ corresponds to a face(t)
  of $B_n$ (under $P$) 
  if and only if it corresponds to a face(t)
  of the representation polytope $P(D)$ (under $D$).
  
  
  
    Let $H \leq S_{n}$ be the stabilizer of a point,
    say $n$. (So $H\iso S_{n-1}$.)
    By the description of the facets of $B_n$, we know that
    $S_n \setminus H = \{g \in S_n\mid g(n)\neq n\}$
    corresponds to a facet of $B_n$.
    Thus $D(S_n\setminus H)$ is a facet of $P(D)$.
      
    Let $\phi$ be the character of $D$.
    The character of the standard permutation representation~$P$
    has the form  $(1_H)^{S_n} = 1_{S_n} + \chi$,
    where $\chi$ is an irreducible character of $S_n$.
    We are going to show that $\chi$ is the only 
    nontrivial irreducible constituent of~$\phi$.
    
    As we remarked in the first paragraph of the proof,
    we can assume that $\phi$ 
    does not contain the trivial character.
    The matrix $\sum_{g\in S_n} D(g)$ is fixed under
    multiplication with elements from $D(S_n)$,
    and since the trivial representation is not a constituent of~$D$,
    we have $\sum_{g\in S_n} D(g) = 0$.
    Geometrically, this means that
    the origin is the barycenter 
    of the representation polytope~$P(D)$.
    As $D(S_n\setminus H)$ is a facet of $P(D)$,
    we must have
      \[ \sum_{g\in S_n \setminus H} D(g) \neq 0
         ,\quad \text{and} \quad 
          \sum_{g\in  H} D(h) \neq 0.
      \]
    It follows that the restricted character
    $\phi_H$ contains the trivial character $1_H$
    as a constituent.
    Using Frobenius reciprocity and the fact that
    $(1_H)^{S_n} = 1_{S_n} + \chi$, we get 
    \[ 0 \neq [\phi_H, 1_H ] = [\phi, (1_H)^{S_n}]
       = [\phi,1_{S_n}] + [\phi,\chi] =  [\phi,\chi].  
    \]
    Thus $\chi$ is a constituent of $\phi$.

  Since dimension is a combinatorial invariant, we must have
  $\dim P(D) = \dim B_n = \chi(1)^2 $.
  On the other hand, we have
  $\dim P(D) = \sum_{\psi} \psi(1)^2$,
  where the sum runs over the nontrivial 
  irreducible constituents $\psi$ of $\phi$,
  not counting multiplicities~\cite[Theorem~3.2]{guralnickperkinson06}.
  It follows that $\chi$ is the only irreducible constituent of $\phi$,
  and thus $D$ and $P$ are stably equivalent.
\end{proof}

\begin{remark}
  In the preceding proof,
  we reduced to the case that the combinatorial isomorphism
  sends $P(g)$ to $D(g)$ (for any $g\in S_n$).
  If we could show that then $P(g)\mapsto D(g)$ 
  can be extended to an affine isomorphism, 
  Lemma~\ref{l:snbirk_unique} would follow
  from a characterization of effective equivalence
  by Baumeister and 
  Grüninger~\cite[Corollary~4.5]{BaumeisterGrueninger15}.
  But we do not know how to do this, or whether this is even true
  more generally 
  (for combinatorial isomorphisms of this form between
  representation polytopes of arbitrary groups).
\end{remark}  



Finally, we prove our main result:

